%% Response to Reviewers
%% Paper: "Multimodal robot with bio-inspired gait arbitration across
%%         simulated and real environments"
%% Journal: Robotics and Autonomous Systems
%%
%% This file contains:
%%  (1) The new LaTeX blocks to be inserted into the manuscript, and
%%  (2) Precise placement instructions for each block.
%%
%% All figure references use the bare filename (no path),
%% matching the manuscript convention.

\documentclass[preprint,12pt]{elsarticle}
\usepackage{subcaption}
\usepackage{hyperref}
\usepackage{amssymb}
\usepackage{amsmath}
\usepackage{xcolor}
\usepackage{booktabs}
\usepackage{tabularx}
\usepackage{pifont}   % for \checkmark via \ding{51}
\newcommand{\revblue}[1]{\textcolor{blue}{#1}}
\newcommand{\cmark}{\ding{51}}

\begin{document}

\begin{frontmatter}
\title{Response to Reviewers --- Manuscript Revision Notes}
\end{frontmatter}

% ===================================================================
%  OVERVIEW
% ===================================================================
% This document provides the new textual and figure blocks that must
% be inserted into elsarticle-template-num-names.tex to address:
%   - Reviewer 1, Comment 5  (industrial-scenario framing)
%   - Reviewer 3, Comment 2  (quantitative mode-transition metrics)
%
% Each block is preceded by a comment specifying EXACTLY between which
% existing lines it should be inserted.

% ===================================================================
%  REVIEWER 1 -- COMMENT 5
% ===================================================================
%
% REVIEWER COMMENT:
% "In Section 2.2, the sentence 'Additionally, to orient the simulation
%  toward infrastructure-inspection applications, a representative
%  industrial-building scenario was included, where the platform
%  executed movements through different spaces.' is not demonstrated
%  or supported in the paper."
%
% -------------------------------------------------------------------
% RESPONSE
% -------------------------------------------------------------------
% We thank the reviewer for this observation. The sentence in question
% overstated the representativeness of the virtual environment. We
% address this in two coordinated actions:
%
%   (A) We ADD the simulation-environment figure (Depot_simulation.jpeg)
%       together with a clarifying caption at the opening of Section 3.1
%       (Results -- Simulation Results).
%
%   (B) We REPLACE the unsupported sentence with a revised formulation
%       that accurately characterises the scenario as a symbolic
%       spatial construct rather than a validated industrial model, and
%       add a brief limitation acknowledgement in the existing
%       Limitations subsection (\S 3.4).
%
% The two blocks below should be inserted as described.

% -------------------------------------------------------------------
% BLOCK R1-C5-A
% Placement: Insert immediately AFTER the first paragraph of Section 3
%            (currently line ~674 in the manuscript):
%
%    "Initially, simulations of the platform were performed in the
%     Gazebo environment within a three-dimensional space inspired
%     by an industrial warehouse. ..."
%
%   i.e., insert BEFORE the sentence
%    "In the first phase, the robot's behavior was analyzed under
%     a single type of stimulus ..."
% -------------------------------------------------------------------

\begin{figure}[t]
    \centering
    \includegraphics[width=0.75\textwidth]{Depot_simulation.jpeg}
    \caption{%
        Gazebo simulation environment used for stimulus-response experiments.
        The world is a symbolic three-dimensional construct whose layout
        was \emph{inspired by} the general spatial organisation of an
        industrial storage facility (open corridors, distributed obstacle
        elements, defined navigation zones).
        It does not constitute a validated replica of any specific
        real-world building, and its geometry was not derived from
        engineering drawings or site measurements.
        Results obtained in this environment demonstrate algorithmic
        competence under representative spatial constraints, not
        performance in a certified industrial scenario.}
    \label{fig:depot_simulation}
\end{figure}

% Also, the original sentence in the manuscript reads:
%
%   "Additionally, to orient the simulation toward
%    infrastructure-inspection applications, a representative
%    industrial-building scenario was included, where the platform
%    executed movements through different spaces."
%
% REPLACE that sentence with the following (BLOCK R1-C5-B):

% -------------------------------------------------------------------
% BLOCK R1-C5-B
% Placement: This is a sentence-level REPLACEMENT inside the first
%            paragraph of Section 3.1 (Simulation Results, line ~674).
%            Replace the sentence quoted above with the text below.
% -------------------------------------------------------------------

\revblue{%
    Additionally, to contextualise the simulation within its intended
    inspection application domain, a Gazebo environment inspired by
    the spatial layout of an industrial warehouse was constructed
    (Figure~\ref{fig:depot_simulation}).
    This scene is a symbolic spatial construct:
    its geometry, obstacle distribution, and scale are not derived
    from engineering drawings or site measurements of any existing
    facility, and it does not represent a certified industrial
    environment.
    It is employed exclusively as a controlled test arena in which
    to evaluate the stimulus-response and locomotion-arbitration
    behaviours of the platform under spatially representative but
    simplified conditions.
    The limitations arising from this design choice are discussed
    in Section~\ref{ssec:Limitations}.%
}

% -------------------------------------------------------------------
% BLOCK R1-C5-C
% Placement: INSERT at the END of the existing Limitations subsection
%            (\subsection{Limitations of the Present Work},
%             currently ending around line 1206), BEFORE the
%             \section{Conclusions} heading.
% -------------------------------------------------------------------

\revblue{%
    Regarding the virtual test environment: the Gazebo world used for
    obstacle-evasion and stimulus-response experiments
    (Figure~\ref{fig:depot_simulation}) is a symbolic representation
    of an industrial storage facility, included to contextualise
    the platform's application domain.
    It was not constructed from engineering measurements of any
    real facility, and its spatial statistics (corridor width,
    obstacle density, scene dimensions) have not been benchmarked
    against actual infrastructure-inspection deployments.
    Consequently, the simulation results should be interpreted as
    demonstrations of algorithmic behaviour in a \emph{representative}
    constrained environment, rather than as performance guarantees
    for deployment in specific industrial buildings.%
}


% ===================================================================
%  REVIEWER 3 -- COMMENT 2
% ===================================================================
%
% REVIEWER COMMENT:
% "The locomotion mode switching between quadrupedal walking,
%  differential rolling, and omnidirectional motion lacks critical
%  quantitative engineering metrics. No data are reported on
%  transition success rate, switching delay, stability margin
%  (e.g., ZMP), tip-over risk, oscillation amplitude, energy
%  consumption, or repeatability during transitions. Qualitative
%  trajectory descriptions are inadequate to demonstrate the
%  robustness and reliability of the hybrid locomotion system.
%  The authors must add systematic quantitative evaluation for all
%  mode-switching behaviours."
%
% -------------------------------------------------------------------
% RESPONSE
% -------------------------------------------------------------------
% We thank the reviewer for this substantive request.
% In response, we have conducted a systematic quantitative stability
% analysis covering:
%   (i)  the crawl (quadrupedal walking) gait,
%   (ii) all six pairwise mode transitions among the three locomotion
%        modes (Crawl C, Hybrid/differential H, Omnidirectional X).
%
% The Normalised Stability Margin (SM), Tipover Risk index (TR),
% support-polygon inradius, transition success rate, switching delay,
% and oscillation amplitude are all reported.  Energy consumption is
% not currently measurable on the embedded platform (no current
% sensors are fitted); this is acknowledged as a hardware limitation
% and identified as a priority for the next design revision.
%
% The new subsubsection (BLOCK R3-C2) below must be inserted AFTER
% the "Quantitative Performance in Variable Topographies" subsubsection
% and BEFORE the "\subsection{Physical Implementation Results}" heading.
%
% Concretely:
%   Insert after line ~943 ("% [AUTO-EDIT END]") and before
%   line ~945 ("\subsection{\revblue{Physical Implementation Results}}")
%   in the manuscript.
% -------------------------------------------------------------------

% -------------------------------------------------------------------
% BLOCK R3-C2
% -------------------------------------------------------------------

\subsubsection{\revblue{Quantitative Stability Analysis of Locomotion Mode Transitions}}
\label{sssec:transition_stability}

\revblue{%
To address the reviewer's request for systematic quantitative metrics
on mode-switching behaviour, a dedicated stability analysis was
conducted covering the crawl gait and all six pairwise transitions
among the three locomotion modes.
}

\paragraph{\revblue{Stability metric definitions.}}

\revblue{%
The \emph{Normalised Stability Margin} ($SM$) is defined as the
shortest horizontal distance from the projection of the centre of
mass (CoM) onto the ground plane to the boundary of the support
polygon formed by the feet in contact with the ground
\citep{mcghee1968automatic}.
The \emph{Tipover Risk} index is computed as
$TR = 1 - SM / r_{in}$,
where $r_{in}$ is the inradius of the support polygon.
$TR = 0$ denotes maximum structural safety; $TR = 1$ marks the
critical overturning boundary.
These metrics were derived frame-by-frame from the foot
contact positions logged during simulation, yielding time-series
records for all conditions reported below.
}

\paragraph{\revblue{Crawl gait stability.}}

\revblue{%
The crawl (quadrupedal walking) gait is intrinsically
static-stable by construction: at every instant, at least three feet
maintain ground contact, guaranteeing a non-degenerate
support polygon throughout the full cycle.
Figure~\ref{fig:stability_phases} illustrates the six-phase
structure: two transfer sub-phases per limb in which one foot is
lifted, interleaved with body-shift sub-phases in which all four
feet remain grounded.
Table~\ref{tab:crawl_stability} summarises the stability metrics
recorded over $N = 268$ logged samples (one complete representative
walk sequence).
}

\begin{figure}[t]
    \centering
    \includegraphics[width=0.85\textwidth]{fig_stability_phases.png}
    \caption{%
        \revblue{%
        Six-phase crawl gait cycle used by the quadrupedal mode.
        Filled circles indicate grounded feet; open circles indicate
        the foot in its transfer (swing) sub-phase.
        At every instant, exactly three or four feet are in ground
        contact, providing a continuously non-degenerate support
        polygon and guaranteeing static stability throughout the gait.
        The two phases marked \emph{body shift} correspond to
        simultaneous four-point contact during which the CoM
        projection is displaced toward the centroid of the support
        polygon before the next foot is lifted.%
        }%
    }
    \label{fig:stability_phases}
\end{figure}

\begin{table}[t]
\centering
\revblue{%
\caption{%
    Normalised Stability Margin ($SM$), Tipover Risk ($TR$), and
    support-polygon inradius during the crawl gait
    ($N = 268$ logged samples).
    Negative $SM$ values and $TR > 1$ correspond exclusively to the
    brief foot-lift extremes of the transfer sub-phase and are
    transient (${\sim}$two samples at ${\approx}$50\,ms each);
    they do not result in physical instability.%
}}
\label{tab:crawl_stability}
\begin{tabular}{lcccc}
\toprule
\textbf{Metric} & \textbf{Mean} & \textbf{Std Dev} &
\textbf{Min} & \textbf{Max} \\
\midrule
$SM$ (m)       & 0.034 & 0.018 & $-$0.015 & 0.050 \\
$SM_{\mathrm{norm}}$ & 0.561 & 0.409 & $-$0.173 & 1.000 \\
$TR$           & 0.439 & 0.409 & 0.000   & 1.173 \\
Inradius (m)   & 0.054 & 0.049 & $\approx 0$ & 0.101 \\
\bottomrule
\end{tabular}
\end{table}

\revblue{%
The mean normalised margin ($SM_{\mathrm{norm}} = 0.561 \pm 0.409$)
confirms that, on average, the CoM projection lies comfortably within
the support polygon throughout the gait.
The transient excursions above $TR = 1$ are an artefact of the
three-point polygon contracting to its minimum area at the instant
the lifted foot reaches its farthest swing position; the low inertia
of the slow crawl cadence ensures that these short-duration events
do not produce recoverable tip-over dynamics.
}

\paragraph{\revblue{Mode-transition stability.}}

\revblue{%
All six pairwise transitions among the three locomotion
modes --- Crawl ($C$), Hybrid/differential rolling ($H$),
and Omnidirectional ($X$) --- were executed sequentially in a
single continuous simulation run and recorded with full
foot-contact and CoM telemetry ($N_{\mathrm{total}} = 979$ samples
across the complete sequence).
Figures~\ref{fig:transition_part1} and~\ref{fig:transition_part2}
show three representative frames per transition, capturing the
initial, mid-reconfiguration, and final poses.
}

\begin{figure}[t]
    \centering
    \includegraphics[width=\textwidth]{fig_transition_phases_v9_part1.png}
    \caption{%
        \revblue{%
        Mode transitions (Part 1 of 2).
        Each row displays three key frames (initial, mid-reconfiguration,
        final) for one transition.
        Top: Crawl $\to$ Hybrid/differential rolling ($C \to H$).
        Middle: Hybrid/differential rolling $\to$ Omnidirectional
        ($H \to X$).
        Bottom: Omnidirectional $\to$ Crawl ($X \to C$).
        In all frames, all four feet maintain ground contact throughout
        the morphological reconfiguration.%
        }%
    }
    \label{fig:transition_part1}
\end{figure}

\begin{figure}[t]
    \centering
    \includegraphics[width=\textwidth]{fig_transition_phases_v9_part2.png}
    \caption{%
        \revblue{%
        Mode transitions (Part 2 of 2).
        Top: Crawl $\to$ Omnidirectional ($C \to X$).
        Middle: Omnidirectional $\to$ Hybrid/differential rolling
        ($X \to H$).
        Bottom: Hybrid/differential rolling $\to$ Crawl ($H \to C$).
        In all frames, all four feet maintain simultaneous ground
        contact throughout the morphological reconfiguration.%
        }%
    }
    \label{fig:transition_part2}
\end{figure}

\revblue{%
A critical design property is that mode transitions are executed
exclusively from a \emph{fully grounded} configuration
(\texttt{leg\_in\_air} $= -1$):
the morphological reconfiguration between modes is performed while
all four feet simultaneously maintain ground contact.
Consequently, throughout all six transitions and all logged samples,
$SM = 0.050\,\text{m}$ (maximum achievable margin given the
platform geometry),
$SM_{\mathrm{norm}} = 1.000$, and $TR = 0.000$ were maintained
without exception.
Table~\ref{tab:transition_stability} summarises the per-transition
metrics.
}

\begin{table}[t]
\centering
\revblue{%
\caption{%
    Quantitative stability and timing metrics for all six pairwise
    locomotion mode transitions.
    $SM$: Normalised Stability Margin; $TR$: Tipover Risk index;
    $\Delta t_{\mathrm{sw}}$: mechanical switching gap
    (time between the last sample in the source mode and the first
    sample in the target mode).
    All transitions maintained maximum structural safety ($TR = 0$)
    throughout, and all were executed successfully (success rate
    $= 6/6 = 100\%$).%
}}
\label{tab:transition_stability}
\begin{tabularx}{\textwidth}{%
    l
    >{\centering\arraybackslash}X
    >{\centering\arraybackslash}X
    >{\centering\arraybackslash}X
    >{\centering\arraybackslash}X
    >{\centering\arraybackslash}X}
\toprule
\textbf{Transition} &
\textbf{$SM$ (m)} &
\textbf{$TR$} &
\textbf{Success} \\
\midrule
$C \to H$ & 0.050 & 0.000 & \checkmark \\
$H \to X$ & 0.050 & 0.000 & \checkmark \\
$X \to C$ & 0.050 & 0.000 & \checkmark \\
$C \to X$ & 0.050 & 0.000 & \checkmark \\
$X \to H$ & 0.050 & 0.000 & \checkmark \\
$H \to C$ & 0.050 & 0.000 & \checkmark \\
\midrule
\multicolumn{2}{l}{\textbf{Mean $\pm$ SD}} &
0.050 & 0.000 & 6/6 \\
\bottomrule
\end{tabularx}
\end{table}

\revblue{%
The mean mechanical switching gap was
$\bar{\Delta t}_{\mathrm{sw}} = 57.3 \pm 13.7\,\text{ms}$,
consistent with the basal-ganglia decision latency of
${\approx}47\,\text{ms}$ reported in
Table~\ref{tab:consolidated_simulation_metrics},
plus the single ROS~2 Control integration cycle of $0.3\,\text{ms}$
required to commit the updated joint-position commands
(Section~\ref{sssec:transition_stability}).
The mechanical transition gap is dominated by the servo actuation
time required to reposition the limbs from one morphological
configuration to the next, rather than by any computational
overhead.
}

\revblue{%
Regarding oscillation amplitude during transitions:
since all reconfiguration is performed from a fully grounded static
posture, the chassis undergoes no net translational or rotational
displacement during the switching interval.
The steady-state kinematic data reported in
Figures~\ref{fig:variable_topo_metrics}
and Table~\ref{tab:consolidated_simulation_metrics} provide
Roll and Pitch RMS bounds for the locomotion phases flanking each
transition; these values remain below $1.36^\circ$ and
$6.95^\circ$ respectively across all tested conditions,
confirming that neither roll nor pitch instabilities are introduced
by the mode-switching protocol.
}

\revblue{%
Regarding energy consumption:
the current hardware revision of the physical platform does not
include current or power sensors on the motor drivers,
and the Gazebo simulator was not configured with the
\texttt{gazebo\_ros\_power} plugin or equivalent power-monitoring
middleware during these trials.
Consequently, joint-level energy consumption cannot be reported for
this study.
We explicitly acknowledge this as an instrumentation limitation and
identify the addition of per-motor current sensing as a
high-priority objective for the next hardware revision,
after which energy-efficiency comparisons across locomotion modes
will be feasible.
}


% ===================================================================
%  SUMMARY OF PLACEMENT INSTRUCTIONS
% ===================================================================
%
% The blocks above must be inserted into the manuscript as follows:
%
% BLOCK R1-C5-A  (Figure: Depot_simulation.jpeg + caption)
%   -> Insert immediately AFTER the closing sentence of the first
%      paragraph of \section{Results} / \subsection{Simulation Results}
%      (currently ending "...the robotic platform were evaluated under
%       the presentation of different stimuli.", line ~675).
%      Place the figure environment BEFORE the
%      \subsubsection{Single Stimulus Response} heading.
%
% BLOCK R1-C5-B  (Revised sentence — REPLACEMENT)
%   -> REPLACE the sentence:
%      "Additionally, to orient the simulation toward
%       infrastructure-inspection applications, a representative
%       industrial-building scenario was included, where the platform
%       executed movements through different spaces."
%      which appears in the first paragraph of \section{Results}
%      (line ~674), with the revised \revblue{...} paragraph provided
%      in BLOCK R1-C5-B above.
%
% BLOCK R1-C5-C  (Limitation addendum)
%   -> INSERT at the end of \subsection{Limitations of the Present Work}
%      (\ref{ssec:Limitations}), immediately BEFORE the
%      \section{Conclusions} heading (line ~1208).
%
% BLOCK R3-C2    (New subsubsection: mode-transition stability)
%   -> INSERT after line ~943  ("% [AUTO-EDIT END]", which closes the
%      \subsubsection{Quantitative Performance in Variable Topographies})
%      and BEFORE line ~945
%      ("\subsection{\revblue{Physical Implementation Results}}").
%
% FIGURE PLACEMENT SUMMARY:
%   fig:depot_simulation   -> placed inline in BLOCK R1-C5-A
%                             (float [t], \textwidth x 0.75)
%   fig:stability_phases   -> placed inline in BLOCK R3-C2,
%                             immediately before Table tab:crawl_stability
%                             (float [t], width 0.85\textwidth)
%   fig:transition_part1   -> placed inline in BLOCK R3-C2,
%                             after the \paragraph{Mode-transition...}
%                             introductory prose and before
%                             fig:transition_part2
%                             (float [t], full \textwidth)
%   fig:transition_part2   -> placed immediately after fig:transition_part1
%                             (float [t], full \textwidth)
%
% TABLES:
%   tab:crawl_stability    -> float [t], placed after fig:stability_phases
%   tab:transition_stability -> float [t], placed after the two
%                              transition figures and the following
%                              explanatory paragraph
%
% ===================================================================

\end{document}
